% ==============================================================================
% The Deliberation Gradient: Mechanistic Signatures of Dual-Process Cognition
% in Large Language Models
%
% ICML 2027 submission
% ==============================================================================

\documentclass[accepted]{article}

% ---------- Packages ----------
\usepackage[utf8]{inputenc}
\usepackage[T1]{fontenc}
\usepackage{amsmath,amssymb,amsfonts}
\usepackage{graphicx}
\usepackage{booktabs}
\usepackage{enumitem}
\usepackage{hyperref}
\usepackage{url}
\usepackage{xcolor}
\usepackage{natbib}
\usepackage{microtype}
\usepackage[margin=1in]{geometry}

% NOTE: Replace \documentclass and geometry with icml2027.sty when available.
% For now we use a standard article class with comparable formatting.
% To switch: \documentclass{icml2027} and remove the geometry package.

% ---------- Custom commands ----------
\newcommand{\sone}{S1-like}
\newcommand{\stwo}{S2-like}
\newcommand{\pzero}{\texttt{P0}}
\newcommand{\ptwo}{\texttt{P2}}
\newcommand{\tend}{\texttt{Tend}}
\newcommand{\todo}[1]{\textcolor{red}{\textbf{[TODO: #1]}}}

% ---------- Title ----------
\title{The Deliberation Gradient: Mechanistic Signatures of\\
  Dual-Process Cognition in Large Language Models}

\author{%
  \textbf{Bright Liu}\textsuperscript{1}\thanks{Equal contribution.} \quad
  \textbf{[Author 2]}\textsuperscript{1}$^*$ \quad
  \textbf{[Author 3]}\textsuperscript{1} \quad
  \textbf{[Author 4]}\textsuperscript{1} \\[2pt]
  \textbf{[Author 5]}\textsuperscript{1} \quad
  \textbf{[Author 6]}\textsuperscript{1} \quad
  \textbf{[Author 7]}\textsuperscript{1} \quad
  \textbf{[Advisor]}\textsuperscript{1,2} \\[6pt]
  \textsuperscript{1}Harvard Undergraduate AI Safety (HUSAI) \quad
  \textsuperscript{2}[Department] \\[2pt]
  \texttt{brightliu@college.harvard.edu}
}

\date{}

\begin{document}

\maketitle

% ============================================================
% Abstract
% ============================================================

\begin{abstract}
Large language models exhibit both heuristic and deliberative processing
patterns on cognitive bias tasks: they fall for intuitive lures on conflict
problems while succeeding on structurally matched controls.  But do these
behavioral patterns reflect genuinely distinct internal computations, or
merely surface-level mimicry?  We present the first unified mechanistic
characterization of dual-process-like cognition in LLMs, combining five
complementary methods---linear probes, sparse autoencoder (SAE) feature
analysis with built-in falsification, attention entropy profiling,
representational geometry, and causal interventions---on a contrastive
cognitive bias benchmark of 284 items across 7 task categories.  Our
headline comparison exploits a natural experiment: Llama-3.1-8B-Instruct
and DeepSeek-R1-Distill-Llama-8B share identical architectures but differ
in reasoning distillation training, isolating the effect of training on
internal representations.  Rather than testing for a binary System~1 versus
System~2 dichotomy, we measure a graded \emph{deliberation-intensity
gradient} across model layers and positions.  [We find that \ldots.
Placeholder for headline results: layer-wise probe accuracy, number of
genuine SAE deliberation features post-falsification, causal steering
effect sizes, and cross-model convergence patterns.]  Our results
[support/complicate/refute] the hypothesis that dual-process-like
computation has identifiable mechanistic substrates in LLMs, with
implications for understanding when models reason versus pattern-match.
\end{abstract}

% ============================================================
% Main text
% ============================================================

% --- Figure 1: Benchmark schematic ---
\begin{figure*}[t]
\centering
% \includegraphics[width=\textwidth]{figures/fig1_benchmark_schematic.pdf}
\fbox{\parbox{0.95\textwidth}{\centering\vspace{3cm}
  \textbf{Figure 1: Benchmark and experimental overview}\\[4pt]
  \small Left: the 7 task categories with example conflict/no-conflict pairs.\\
  Center: activation extraction at key token positions (P0, P2, T0--Tend).\\
  Right: five analysis workstreams applied to the shared activation cache.
\vspace{3cm}}}
\caption{Overview of the experimental framework.  (a)~The contrastive
  benchmark presents matched conflict/no-conflict problem pairs across 7
  cognitive bias categories.  (b)~Residual stream activations are extracted
  at semantically meaningful token positions for all four models.
  (c)~Five complementary mechanistic methods probe the deliberation-intensity
  gradient from different angles.}
\label{fig:overview}
\end{figure*}

% ==============================================================================
% Introduction
% ==============================================================================

\section{Introduction}
\label{sec:introduction}

Humans solve problems using qualitatively different modes of processing.
Kahneman's influential dual-process framework distinguishes a fast, effortless,
heuristic mode (``System~1'') from a slow, effortful, deliberative mode
(``System~2'') \citep{kahneman2011thinking}.  While the binary framing has
faced sustained criticism---Melnikoff and Bargh argue that the ``mythical number
two'' obscures what is better understood as a continuum of processing intensity
\citep{melnikoff2018mythical}---the core observation remains: cognitive systems
flexibly allocate effort, and this allocation has measurable consequences for
accuracy, latency, and susceptibility to systematic biases
\citep{evans2013dual, deneys2012bias}.

Large language models (LLMs) now exhibit strikingly analogous behavioral
patterns.  On classic cognitive bias tasks, instruction-tuned models reproduce
human-like errors on conflict items---problems where a salient heuristic cue
diverges from the normatively correct answer---while succeeding on structurally
matched controls where heuristic and deliberative cues converge
\citep{hagendorff2023human}.  Reasoning-trained models, such as those distilled
from DeepSeek-R1 \citep{deepseekr1}, resist many of these heuristic lures,
producing explicit chain-of-thought (CoT) traces that superficially resemble
deliberative processing.  A growing body of work has documented these behavioral
parallels across models and task families \citep{xie2025cogsci, rao2024cobbler,
codaforno2025dual, ziabari2025reasoning}.

Yet behavior alone cannot distinguish genuine processing-mode differences from
surface mimicry.  A model that memorizes the correct answer to the bat-and-ball
problem produces the same output as one that ``reasons through'' it, but via
radically different internal computation.  Moreover, chain-of-thought traces are
often unfaithful to the underlying computation: recent work shows that the
majority of reasoning steps in CoT---including apparent self-correction
markers---do not causally influence the final answer
\citep{lanham2023measuring, jiang2025aha}.  The critical open question is
therefore not whether LLMs \emph{behave} as if they have fast and slow
processing modes, but whether such modes have \emph{mechanistic correlates} in
the model's internal representations.

Two recent lines of work begin to address this gap.  On the mechanistic side,
probing studies have shown that LLM hidden states encode self-verification
signals before the model generates its answer \citep{zhang2025probing}, that
attention heads specialize for recall versus reasoning
\citep{fartale2025recall}, and that sparse autoencoder (SAE) features can
identify and steer reasoning-related computations
\citep{fang2026sae, yang2026step, goodfire2025r1}.  On the cognitive science
side, CogBias \citep{huang2026cogbias} recently combined linear probes with
activation steering to detect and mitigate cognitive biases in LLM internal
representations.  However, no prior work has attempted a \emph{unified}
mechanistic characterization of dual-process-like computation, combining
multiple complementary interpretability methods on a benchmark explicitly
designed to elicit and control for the distinction.

This paper presents the first such characterization.  We make four
contributions:

\begin{enumerate}[leftmargin=*,itemsep=2pt]
  \item \textbf{A contrastive cognitive bias benchmark} of 284 items (142
  conflict/no-conflict matched pairs) across 7 task categories, with novel
  structural isomorphs designed to resist memorization-based shortcuts
  (Section~\ref{sec:benchmark}).

  \item \textbf{A multi-method mechanistic analysis} combining five
  complementary workstreams---linear probes, SAE feature analysis, attention
  entropy, representational geometry, and causal interventions---applied to four
  models spanning standard and reasoning-trained architectures
  (Section~\ref{sec:methods}).

  \item \textbf{A controlled architecture comparison} between
  Llama-3.1-8B-Instruct and DeepSeek-R1-Distill-Llama-8B, which share
  identical architectures (32 layers, 32 attention heads, 4096 hidden
  dimensions) but differ in training.  Any internal differences are attributable
  to reasoning distillation, not architectural confounds.

  \item \textbf{Built-in falsification} at every stage: Hewitt--Liang control
  tasks for probes \citep{hewitt2019control}, the Ma et al. framework for SAE
  feature validation \citep{ma2026falsification}, random-projection baselines
  for geometry, and random-direction controls for causal interventions.
\end{enumerate}

Critically, we frame our investigation not as testing for a binary System~1
versus System~2 dichotomy, but as measuring the \textbf{deliberation-intensity
gradient}: the degree to which models exhibit mechanistically distinct internal
signatures when processing tasks that differentially recruit heuristic versus
deliberative computation.  This graded framing is both more scientifically
defensible and more informative---it asks not ``does a model have System~2?''
but ``how, where, and when does deliberation intensity vary across the network's
representations?''

% ==============================================================================
% Related Work
% ==============================================================================

\section{Related Work}
\label{sec:related_work}

\paragraph{Dual-process theory and LLMs.}
The application of dual-process theory to artificial systems has generated
considerable recent interest.  \citet{lampinen2025nature} provide a
comprehensive review of the connections between dual-process frameworks and
modern AI, arguing that the distinction is better understood as a continuum of
processing characteristics rather than a categorical divide.
\citet{hagendorff2023human} demonstrated that GPT-3.5 reproduces human-like
performance on the Cognitive Reflection Test (CRT), including susceptibility to
heuristic lures, while GPT-4 partially overcomes these biases---a pattern
reminiscent of individual differences in human ``cognitive reflection.''
\citet{codaforno2025dual} tested whether prompting strategies that differentially
recruit fast versus slow processing produce distinct internal representations,
finding overlapping but partially separable subspaces.
\citet{ziabari2025reasoning} showed that interpolating between a base model and
a reasoning-fine-tuned model yields a monotonic gradient of behavioral change on
dual-process tasks, supporting a continuous rather than binary characterization.
Our work moves from behavioral characterization to mechanistic analysis,
testing whether these behavioral gradients have identifiable internal
correlates.

\paragraph{Probing LLM internal states.}
Linear probing has become a standard tool for studying what information is
encoded in neural network representations
\citep{hewitt2019control, burns2022discovering}.  Several recent studies probe
representations specifically related to reasoning and self-monitoring.
\citet{zhang2025probing} showed that hidden states in DeepSeek-R1 encode
self-verification signals---the model's internal state ``knows'' whether its
answer is correct before producing it.
\citet{afzal2025knowing} demonstrated that models internally prepare for
difficult outputs several tokens in advance, suggesting a form of prospective
processing.
\citet{fartale2025recall} disentangled attention heads specialized for factual
recall from those involved in multi-step reasoning.
Our probing analysis extends this line of work to the S1/S2 distinction
specifically, asking whether the conflict versus no-conflict condition is
linearly decodable from residual stream activations and whether this
decodability varies systematically across layers and models.

\paragraph{SAE-based analysis of reasoning.}
Sparse autoencoders (SAEs) decompose neural network activations into
interpretable features, offering finer-grained access to model computation
than linear probes.
\citet{goodfire2025r1} identified specific SAE features in DeepSeek-R1
associated with backtracking, self-reference, and answer urgency, providing
the first mechanistic decomposition of chain-of-thought reasoning.
\citet{fang2026sae} demonstrated that steering SAE features can directly
modulate reasoning quality, while \citet{yang2026step} analyzed how SAE
feature activations evolve across reasoning steps.
Critically, \citet{ma2026falsification} showed that 45--90\% of claimed
``reasoning features'' in SAEs are spurious---they activate in response to
specific tokens rather than the reasoning process itself.  We adopt their
falsification framework as a mandatory control throughout our SAE analysis.

\paragraph{Cognitive biases as interpretability benchmarks.}
Cognitive bias tasks are particularly well-suited to mechanistic interpretability
because they provide \emph{contrastive} conditions: the same task structure can
be presented with or without a heuristic lure, enabling clean comparison of
internal processing on matched stimuli.
\citet{rao2024cobbler} benchmarked cognitive biases in LLMs-as-evaluators.
Most relevant to our work, \citet{huang2026cogbias} recently combined linear
probes with activation steering on a cognitive bias benchmark, achieving
26--32\% bias reduction via internal intervention.  Our work differs in three
respects: (i)~we study the nature of the S1/S2-like processing distinction
rather than steering models away from biases; (ii)~we employ five
complementary mechanistic methods rather than probes alone; and (iii)~we
systematically compare standard and reasoning-trained models sharing the
same architecture.

\paragraph{Chain-of-thought faithfulness.}
The reliability of CoT reasoning as a window into model computation is a
prerequisite for interpreting thinking traces.
\citet{lanham2023measuring} found that many CoT steps can be removed without
affecting the final answer, suggesting low faithfulness.
\citet{jiang2025aha} showed that only ${\sim}2.3$\% of reasoning steps
causally influence model outputs, and that apparent self-correction markers
(``wait,'' ``actually'') are largely decorative.
These findings motivate our emphasis on \emph{activation-level} rather than
\emph{token-level} analysis: we extract representations from the residual
stream and attention patterns, not from the content of the generated text.

% ==============================================================================
% Benchmark
% ==============================================================================

\section{Benchmark}
\label{sec:benchmark}

We construct a \emph{contrastive} cognitive bias benchmark designed to elicit
and control for the S1-like/S2-like processing distinction.  The design follows
three principles: (1)~every conflict item has a structurally matched no-conflict
control \citep{deneys2012bias}; (2)~all items use novel surface features to
resist memorization-based shortcuts \citep{hagendorff2023human}; and (3)~a $2
\times 2$ difficulty $\times$ congruence design disentangles processing mode
from task difficulty.

\subsection{Task Categories}

The benchmark spans seven categories, each targeting a well-characterized
heuristic-deliberative dissociation from the cognitive science literature:

\begin{enumerate}[leftmargin=*,itemsep=1pt]
  \item \textbf{CRT variants} (48 conflict / 48 no-conflict).  Structural
  isomorphs of the Cognitive Reflection Test with modified surface features
  (changed objects, domains, and numerical values).  We draw on the 150
  decontaminated items from \citet{hagendorff2023human}'s OSF repository and
  generate additional structural variants via controlled template perturbation.

  \item \textbf{Base rate neglect} (22 / 22).  Bayesian reasoning problems
  where a diagnostic description cues a low-base-rate category.  Conflict
  items set the base rate against the description; no-conflict items align
  them.  All posterior probabilities are analytically verified.

  \item \textbf{Belief-bias syllogisms} (20 / 20).  Syllogisms crossing
  logical validity ($\times 2$) with conclusion believability ($\times 2$),
  following the Evans taxonomy.  Conflict items have valid--unbelievable or
  invalid--believable conclusions; controls have congruent
  validity--believability.

  \item \textbf{Anchoring} (14 / 14).  Numerical estimation tasks preceded by
  a high or low anchor.  Conflict items use extreme anchors designed to
  distort estimates; no-conflict items use anchors near the correct value.

  \item \textbf{Framing effects} (14 / 14).  Gain- versus loss-framed
  decision problems.  Conflict items present logically equivalent options in
  different frames; controls present options where framing and value align.

  \item \textbf{Conjunction fallacy} (12 / 12).  Probability judgment tasks
  where a detailed description makes the conjunction of two events seem more
  likely than a single constituent.  Novel scenarios replace the classic
  ``Linda problem.''

  \item \textbf{Multi-step arithmetic} (12 / 12).  Graded-difficulty
  arithmetic (1--5 operations) with and without carrying traps.  Conflict
  items embed a salient but incorrect shortcut; no-conflict items do not.
\end{enumerate}

\subsection{Benchmark Statistics}

The final benchmark comprises \textbf{284 items} organized as \textbf{142
matched pairs} (conflict + no-conflict).  Each item includes 2--5
paraphrases for a total of \textbf{558 unique prompts}, allowing us to
assess robustness to surface variation.  Items are presented in two formats:
constrained free-response (``Answer with just a number:'') for activation
analysis, and multiple-choice for behavioral validation.

\paragraph{Scoring.}  We use a three-tier scoring system: (1)~\emph{binary}
(correct/incorrect); (2)~\emph{categorical} (correct / S1-lure / other-wrong
/ refusal)---the primary target for probing; and (3)~\emph{continuous}
(distance from the S2-correct vs.\ S1-predicted answer).

\paragraph{Contamination controls.}  We include classic problems (e.g., the
original bat-and-ball) as contamination baselines only---they are excluded from
all primary analyses.  We test memorization by prompting models to complete the
original CRT verbatim and comparing performance on classic versus novel variants.

\subsection{Example Items}

Table~\ref{tab:examples} illustrates the conflict/no-conflict matched-pair
structure for three representative categories.

\begin{table}[t]
\centering
\caption{Example conflict/no-conflict item pairs from three benchmark
  categories.  Each pair shares the same correct answer; they differ only in
  whether a heuristic lure is present.}
\label{tab:examples}
\vspace{4pt}
\small
\begin{tabular}{p{1.2cm}p{3.0cm}p{3.0cm}cc}
\toprule
\textbf{Category} & \textbf{Conflict item} & \textbf{No-conflict control}
  & \textbf{Lure} & \textbf{Correct} \\
\midrule
CRT &
  A notebook and a pen cost \$2.20 together.  The notebook costs \$2.00
  more than the pen.  How much does the pen cost? &
  A notebook and a pen cost \$2.10 together.  The notebook costs \$2.00.
  How much does the pen cost? &
  \$0.20 & \$0.10 \\[6pt]
Base rate &
  In a group of 1000 people, 5 are engineers and 995 are lawyers.  Tom
  likes puzzles and circuits.  Is Tom more likely an engineer or a
  lawyer? &
  In a group of 1000 people, 500 are engineers and 500 are lawyers.  Tom
  likes puzzles and circuits.  Is Tom more likely an engineer or a
  lawyer? &
  Engineer & Lawyer \\[6pt]
Syllogism &
  All flowers need water.  Roses need water.  Therefore roses are
  flowers.  Is this valid? &
  All mammals are warm-blooded.  Dogs are mammals.  Therefore dogs are
  warm-blooded.  Is this valid? &
  Valid & Invalid \\
\bottomrule
\end{tabular}
\end{table}

\subsection{Scoring Detail}

We employ three scoring tiers, applied automatically via regex patterns
stored alongside each benchmark item:
\begin{enumerate}[leftmargin=*,itemsep=1pt]
  \item \textbf{Binary}: 1 if the model's response matches the
  \texttt{correct\_answer} regex, 0 otherwise.
  \item \textbf{Categorical}: one of \{correct, S1-lure, other-wrong,
  refusal\}.  The S1-lure label is assigned when the response matches the
  \texttt{lure\_answer} regex (conflict items only).  Refusals are detected
  via keyword matching.  This four-way classification is the primary probing
  target.
  \item \textbf{Continuous} (numerical categories only):
  $|r - c| / |l - c|$, where $r$ is the extracted numerical response, $c$
  is the correct answer, and $l$ is the lure answer.  Values near 0
  indicate correctness; values near 1 indicate the lure response; values
  $> 1$ indicate errors worse than the lure.
\end{enumerate}

\subsection{Behavioral Validation}

\textbf{[Placeholder: behavioral results on all four models.]}  We report the
proportion of S1-lure responses on conflict versus no-conflict items, broken
down by category and model.  The benchmark is validated if (i)~models produce
$>30$\% lure responses on conflict items and (ii)~standard models show
significantly higher lure rates than reasoning-trained models.
% Expected: Table 1 with behavioral accuracy by model x category x condition

% ==============================================================================
% Methods
% ==============================================================================

\section{Methods}
\label{sec:methods}

\subsection{Models}
\label{sec:models}

We study four models spanning two architecture families and two training
regimes (Table~\ref{tab:models}).  The headline comparison is between
Llama-3.1-8B-Instruct \citep{llama31} and DeepSeek-R1-Distill-Llama-8B
\citep{deepseekr1}: identical architectures (32 transformer layers, 32
query heads with 8 KV heads via grouped-query attention, 4096 hidden
dimensions) differing only in that the latter was distilled from a
reasoning-trained teacher.  Any representational differences between these
two models are attributable to reasoning distillation training, not
architectural confounds.  We include Gemma-2-9B-IT \citep{gemma2} and
R1-Distill-Qwen-7B \citep{deepseekr1} as cross-architecture controls.

\begin{table}[t]
\centering
\caption{Models used in this study.  The Llama pair (rows 1 and 3) shares
  identical architecture, enabling controlled comparison.}
\label{tab:models}
\vspace{4pt}
\small
\begin{tabular}{lccccc}
\toprule
\textbf{Model} & \textbf{Type} & \textbf{Layers} & \textbf{Heads} & \textbf{Hidden} & \textbf{SAE source} \\
\midrule
Llama-3.1-8B-Instruct & Standard & 32 & 32 & 4096 & Llama Scope \\
Gemma-2-9B-IT & Standard & 42 & 16 & 3584 & Gemma Scope \\
R1-Distill-Llama-8B & Reasoning & 32 & 32 & 4096 & Llama Scope\textsuperscript{$\dagger$} \\
R1-Distill-Qwen-7B & Reasoning & 28 & 28 & 3584 & Custom \\
\bottomrule
\end{tabular}
\vspace{2pt}
\raggedright
\footnotesize{\textsuperscript{$\dagger$}Llama Scope SAEs were trained on the
base model; we verify reconstruction fidelity on the distilled model before
use.}
\end{table}

\subsection{Activation Extraction}
\label{sec:extraction}

For each model--problem pair, we extract residual stream activations (the
output of each transformer block, post-MLP and post-residual connection) at
all layers and at a set of semantically meaningful token positions.  For
standard models, we extract at two positions: \textbf{P0} (last token of the
prompt, before generation begins) and \textbf{P2} (the answer token).
For reasoning models, we additionally extract at positions within the
thinking trace: \textbf{T0} (first token after \texttt{<think>}),
\textbf{T25/T50/T75} (25\%/50\%/75\% through the thinking trace), and
\textbf{Tend} (last token before \texttt{</think>}).

Activations are stored in BFloat16 in a shared HDF5 cache following a
fixed schema, along with per-head attention entropy metrics computed
incrementally during generation (we never materialize full attention
matrices).  The complete schema is specified in our data contract
documentation.  Total storage for the full benchmark across all four models
is approximately 4\,GB.

\subsection{Workstream 1: Linear Probes}
\label{sec:probes}

We train $\ell_2$-regularized logistic regression probes
(\texttt{LogisticRegressionCV} with 5-fold stratified cross-validation,
stratified by task category and target label) on residual stream activations
at each layer and token position.  We probe for four targets:

\begin{enumerate}[leftmargin=*,itemsep=1pt]
  \item \textbf{Task type}: conflict vs.\ no-conflict condition.
  \item \textbf{Correctness}: whether the model will answer correctly.
  \item \textbf{Bias susceptibility}: whether the model will produce the
  specific S1-lure answer (on conflict items only).
  \item \textbf{Processing mode}: within matched pairs, distinguishing the
  conflict from the no-conflict version of structurally identical problems.
\end{enumerate}

\paragraph{Controls.}  We apply three mandatory controls:
(i)~\emph{Hewitt--Liang control tasks} \citep{hewitt2019control}: probes
trained on randomly permuted labels establish a floor for probe
expressiveness.  Selectivity (true accuracy minus control accuracy) below
5 percentage points indicates the signal reflects probe capacity, not
representational structure.
(ii)~\emph{Permutation tests}: 1000 label shuffles per layer to establish
significance thresholds, using the North et al.\ correction
\citep{north2002note}.
(iii)~\emph{Leave-one-category-out}: training on 6 of 7 task categories and
testing on the held-out category to assess whether probes learn a
domain-general S1/S2 signal rather than category-specific features.

We report ROC-AUC as the primary metric with bootstrap 95\% confidence
intervals (1000 resamples).  We additionally train 2-layer MLP probes
(hidden dimension 256, ReLU activation) as a nonlinearity check and apply
Contrast-Consistent Search \citep[CCS;][]{burns2022discovering} as an
unsupervised corroboration for the correctness target.

\subsection{Workstream 2: SAE Feature Analysis}
\label{sec:sae}

We use pre-trained sparse autoencoders (Llama Scope for Llama-based models
\citep{he2024llama}; Gemma Scope for Gemma \citep{lieberum2024gemma}) to
decompose residual stream activations into sparse feature vectors.  For each
SAE feature, we compute the mean activation on conflict versus no-conflict
items and test for differential activation via Mann--Whitney~$U$ tests,
corrected for multiple comparisons using the Benjamini--Hochberg procedure
at $q = 0.05$ \citep{benjamini1995controlling}.

\paragraph{Visualization.}  We construct volcano plots (log fold-change vs.\
$-\log_{10} p$) to identify features with both large effect sizes and
statistical significance.  Features in the upper corners---large differential
activation between conditions---are candidate ``deliberation features'' and
``heuristic features.''

\paragraph{Falsification.}  Following \citet{ma2026falsification}, who
showed that 45--90\% of claimed SAE reasoning features are spurious
token-level artifacts, we apply their falsification protocol to every
candidate feature.  For each feature, we identify its top-3 activating tokens,
inject those tokens into 100 random non-cognitive-bias texts, and test whether
the feature still activates.  Features that survive this filter are classified
as genuine processing-mode features; we report falsification rates for all
candidates.

\subsection{Workstream 3: Attention Entropy}
\label{sec:attention}

We quantify attention focus via Shannon entropy
$H_{l,h,t} = -\sum_i a_i \log_2 a_i$ and the Gini coefficient of attention
weights, computed per-head per-position during generation.  We normalize
entropy by $\log_2 t$ (where $t$ is the current sequence length) to control
for the mechanical effect of longer sequences on entropy.

For each head, we test whether the attention metric distribution differs
between conflict and no-conflict items using Mann--Whitney~$U$ tests,
correcting across all heads $\times$ metrics via BH-FDR.  We classify heads
as \emph{S2-specialized} if they show significantly higher entropy on conflict
items in $\geq 3$ of 5 metrics with effect size $|r_{\text{rb}}| \geq 0.3$,
and \emph{S1-specialized} for the opposite pattern.

\paragraph{GQA handling.}  Because Llama uses grouped-query attention (32
query heads sharing 8 KV heads in groups of~4), heads within the same KV
group are not statistically independent.  We report results at both the
query-head and KV-group granularity, with the latter being the conservative
estimate.

\paragraph{Gemma-2 sliding window.}  Gemma-2's odd-numbered layers use a
4096-token sliding window for attention.  We analyze odd and even layers
separately and never pool them.

\subsection{Workstream 4: Representational Geometry}
\label{sec:geometry}

We analyze the geometric structure of activation point clouds corresponding to
S1-like and S2-like processing conditions.

\paragraph{Clustering.}  At each layer, we compute the \emph{cosine
silhouette score} between conflict and no-conflict activation clusters.  We
first reduce dimensionality to 50--100 principal components to avoid the
$d \gg N$ pitfall (with $d = 4096$ and $N \approx 284$, any two random
classes are linearly separable by Cover's theorem).  All silhouette scores
are compared against a permutation baseline (10,000 label shuffles).

\paragraph{Cross-model comparison.}  We compute layer-matched Centered Kernel
Alignment \citep[CKA;][]{kornblith2019similarity} between
Llama-3.1-8B-Instruct and R1-Distill-Llama-8B on conflict and no-conflict
items separately.  We predict that CKA diverges more between the two models
on conflict items (where reasoning distillation should produce the largest
representational changes) than on no-conflict items.

\paragraph{Intrinsic dimensionality.}  We estimate the intrinsic
dimensionality of S1-like and S2-like activation point clouds using the
Two-NN estimator \citep{facco2017estimating}.  Lower intrinsic
dimensionality for S2-like activations would suggest a more stereotyped
deliberation circuit; higher dimensionality would suggest diverse processing
strategies.

\paragraph{Random projection baseline.}  For any UMAP visualization, we
verify that apparent cluster structure survives projection to 2D via random
Gaussian matrices (100 draws).  If clusters appear in random projections,
the structure is genuine; otherwise, UMAP may be hallucinating separation.

\subsection{Workstream 5: Causal Interventions}
\label{sec:causal}

To move from correlation to causation, we intervene on SAE features and
measure behavioral effects.

\paragraph{Feature steering.}  For each candidate deliberation feature
identified in Workstream~2 (post-falsification), we perform activation
steering: we add $\alpha \cdot \hat{f}$ to the residual stream at the
feature's peak layer during inference, where $\hat{f}$ is the SAE decoder
direction and $\alpha$ is a steering coefficient.  We vary $\alpha$ over
$\{0.5, 1.0, 2.0, 3.0, 5.0\}$ to obtain dose-response curves.

\paragraph{Controls.}  We compare against (i)~steering along a randomly
sampled direction in activation space (same norm), which should produce
negligible behavioral change if the identified features are causally
specific; and (ii)~ablating the candidate features (setting their
activation to zero), which should shift behavior in the opposite direction
from amplification.

\paragraph{Metrics.}  We measure the change in (i)~probability of the
correct answer, (ii)~probability of the S1-lure answer, and (iii)~general
capability on a held-out set of non-bias tasks (to check for side effects).
All experiments use 3 random seeds.

% ==============================================================================
% Results
% ==============================================================================

\section{Results}
\label{sec:results}

% -------------------------------------------------------------------
\subsection{Behavioral Validation}
\label{sec:results_behavioral}

\textbf{[Result placeholder.]}  Table~\ref{tab:behavioral} reports model
accuracy and S1-lure rates across all seven task categories.

\begin{table}[t]
\centering
\caption{Behavioral results: accuracy (\%) and S1-lure rate (\%) on conflict
  and no-conflict items.  Standard models are expected to show higher lure
  rates on conflict items; reasoning models should partially overcome this
  pattern.}
\label{tab:behavioral}
\vspace{4pt}
\small
\begin{tabular}{lcccc}
\toprule
& \multicolumn{2}{c}{\textbf{Conflict}} & \multicolumn{2}{c}{\textbf{No-Conflict}} \\
\cmidrule(lr){2-3}\cmidrule(lr){4-5}
\textbf{Model} & Acc. & Lure \% & Acc. & Lure \% \\
\midrule
Llama-3.1-8B-Instruct   & -- & -- & -- & -- \\
Gemma-2-9B-IT            & -- & -- & -- & -- \\
R1-Distill-Llama-8B     & -- & -- & -- & -- \\
R1-Distill-Qwen-7B      & -- & -- & -- & -- \\
\bottomrule
\end{tabular}
\end{table}

% -------------------------------------------------------------------
\subsection{Linear Probes}
\label{sec:results_probes}

\textbf{[Result placeholder.]}  Figure~\ref{fig:probe_accuracy} shows
layer-wise probe ROC-AUC for the conflict/no-conflict classification target
across all four models.

\begin{figure}[t]
\centering
% \includegraphics[width=\linewidth]{figures/fig2_probe_accuracy.pdf}
\fbox{\parbox{0.95\linewidth}{\centering\vspace{2cm}
  \textbf{Figure 2: Layer-wise probe accuracy}\\[4pt]
  \small X-axis: layer index.  Y-axis: ROC-AUC.\\
  4 model curves with bootstrap 95\% CIs.\\
  Dashed line: Hewitt--Liang control task performance.\\
  Gray band: permutation test 95th percentile.
\vspace{2cm}}}
\caption{Layer-wise probe ROC-AUC for classifying conflict versus
  no-conflict items from residual stream activations at position~P0.
  Shaded bands show bootstrap 95\% CIs.  The dashed line indicates
  Hewitt--Liang control task performance (selectivity floor).  [We find
  that \ldots]}
\label{fig:probe_accuracy}
\end{figure}

\paragraph{Selectivity.}  [Report selectivity = real accuracy $-$ control
accuracy at the peak layer for each model.  Only selectivity $>$ 5pp is
meaningful.]

\paragraph{Cross-category generalization.}  [Report leave-one-category-out
accuracy.  If it holds, the probe learned a domain-general signal.]

\paragraph{Nonlinearity check.}  [Compare logistic regression vs.\ MLP probe
accuracy.  If MLP substantially outperforms, the signal has a nonlinear
component.]

% -------------------------------------------------------------------
\subsection{SAE Feature Analysis}
\label{sec:results_sae}

\textbf{[Result placeholder.]}  Figure~\ref{fig:volcano} shows the volcano
plot of differential SAE feature activation.

\begin{figure}[t]
\centering
% \includegraphics[width=\linewidth]{figures/fig3_sae_volcano.pdf}
\fbox{\parbox{0.95\linewidth}{\centering\vspace{2cm}
  \textbf{Figure 3: SAE volcano plot}\\[4pt]
  \small X-axis: log fold-change (conflict vs.\ no-conflict).\\
  Y-axis: $-\log_{10}(p)$ (BH-FDR corrected).\\
  Upper-left: ``heuristic features.''\\
  Upper-right: ``deliberation features.''\\
  Red points: features surviving Ma et al.\ falsification.\\
  Gray points: features failing falsification.
\vspace{2cm}}}
\caption{Volcano plot of SAE feature activation differences between
  conflict and no-conflict items (Llama-3.1-8B-Instruct, layer~[X]).
  Features in the upper corners show large, statistically significant
  differential activation.  Red: features surviving falsification
  \citep{ma2026falsification}; gray: spurious features eliminated by the
  falsification protocol.  [We identify \ldots]}
\label{fig:volcano}
\end{figure}

\paragraph{Falsification rates.}  [Report: of $N$ candidate features, $M$
($X$\%) survived the Ma et al.\ falsification protocol.  Compare to Ma et
al.'s reported 45--90\% spurious rate.]

\paragraph{Feature interpretation.}  [For the top-$K$ surviving features,
report (i)~top activating examples, (ii)~decoder direction cosine similarity
to known features, (iii)~activation patterns across task categories.]

% -------------------------------------------------------------------
\subsection{Causal Interventions}
\label{sec:results_causal}

\textbf{[Result placeholder.]}  Figure~\ref{fig:causal} shows the effect
of amplifying and ablating candidate deliberation features on model behavior.

\begin{figure}[t]
\centering
% \includegraphics[width=\linewidth]{figures/fig4_causal_bars.pdf}
\fbox{\parbox{0.95\linewidth}{\centering\vspace{2cm}
  \textbf{Figure 4: Causal intervention results}\\[4pt]
  \small Grouped bar chart.\\
  Groups: Amplify S2 features / Ablate S2 features / Random direction.\\
  Bars: $\Delta$ P(correct) and $\Delta$ P(lure).\\
  Error bars: $\pm$ 1 SE across 3 seeds.
\vspace{2cm}}}
\caption{Causal intervention results.  Amplifying candidate deliberation
  features [increases/does not change] $P(\text{correct})$ on conflict
  items, while ablating them [increases/does not change]
  $P(\text{S1-lure})$.  Steering along a random direction of the same norm
  produces negligible change ($\Delta < [X]$).}
\label{fig:causal}
\end{figure}

\paragraph{Dose-response.}  [Report: $\Delta P(\text{correct})$ as a
function of steering coefficient $\alpha$.  Expected: monotonic increase
to a plateau, then degradation at extreme values.]

\paragraph{Side effects.}  [Report: change in general capability on a
held-out task set.  Steering should not degrade non-bias performance.]

% -------------------------------------------------------------------
\subsection{Attention Entropy}
\label{sec:results_attention}

\textbf{[Result placeholder.]}  Figure~\ref{fig:entropy} shows the
distribution of per-head attention entropy on conflict versus no-conflict
items.

\begin{figure}[t]
\centering
% \includegraphics[width=\linewidth]{figures/fig5_attention_entropy.pdf}
\fbox{\parbox{0.95\linewidth}{\centering\vspace{2cm}
  \textbf{Figure 5: Attention entropy ridge plots}\\[4pt]
  \small Y-axis: layer groups (early / mid / late).\\
  X-axis: normalized attention entropy.\\
  Blue: conflict items.  Orange: no-conflict items.\\
  One panel per model.
\vspace{2cm}}}
\caption{Distribution of normalized attention entropy on conflict (blue) and
  no-conflict (orange) items, grouped by layer position.  [We find that
  \ldots]}
\label{fig:entropy}
\end{figure}

\paragraph{Head classification.}  [Report: proportion of heads classified as
S2-specialized, S1-specialized, or unspecialized for each model.  Compare
at query-head and KV-group granularity.]

\paragraph{Temporal dynamics.}  [For reasoning models: compare entropy
trajectories within the thinking trace on conflict vs.\ no-conflict items.]

% -------------------------------------------------------------------
\subsection{Representational Geometry}
\label{sec:results_geometry}

\textbf{[Result placeholder.]}  Figure~\ref{fig:geometry} shows the
layer-wise cosine silhouette score and PCA visualizations.

\begin{figure}[t]
\centering
% \includegraphics[width=\linewidth]{figures/fig6_geometry.pdf}
\fbox{\parbox{0.95\linewidth}{\centering\vspace{2cm}
  \textbf{Figure 6: Representational geometry}\\[4pt]
  \small Left: layer-wise cosine silhouette score (4 models).\\
  Right: PCA visualization at peak-separation layer.\\
  Permutation baseline shown as dashed line.
\vspace{2cm}}}
\caption{Left: layer-wise cosine silhouette score for conflict vs.\
  no-conflict activation clusters (after PCA to 50 dimensions), with
  permutation baseline (dashed).  Right: first two principal components
  at the layer of maximum separation for Llama-3.1-8B-Instruct.  [We find
  that \ldots]}
\label{fig:geometry}
\end{figure}

\paragraph{CKA.}  [Report: layer-matched CKA between Llama-3.1-8B-Instruct
and R1-Distill-Llama-8B, separately for conflict and no-conflict items.]

\paragraph{Intrinsic dimensionality.}  [Report: Two-NN estimates for S1-like
vs.\ S2-like activation clouds at the peak-separation layer.]

% -------------------------------------------------------------------
\subsection{Cross-Workstream Convergence}
\label{sec:results_convergence}

\textbf{[Result placeholder.]}  [Synthesize: which layers show consistent
S1/S2 signals across probes, SAE features, attention entropy, and geometry?
Do the methods agree on which layers and positions carry the most
deliberation-intensity information?  Do they agree across models?]

% ==============================================================================
% Discussion
% ==============================================================================

\section{Discussion}
\label{sec:discussion}

\textbf{[Placeholder: to be completed after results are available.  The
following is the planned structure.]}

\paragraph{The deliberation gradient.}
[Discuss: do the results support a graded deliberation-intensity dimension
rather than a binary S1/S2 dichotomy?  Is the gradient continuous across
layers?  Does it vary smoothly with task difficulty or show sharp transitions?
Connect to \citet{melnikoff2018mythical}'s critique of the binary framework
and \citet{ziabari2025reasoning}'s ``Reasoning on a Spectrum'' finding.]

\paragraph{Standard vs.\ reasoning-trained models.}
[Discuss: how does reasoning distillation change the deliberation gradient?
Does R1-Distill-Llama-8B show (a) a stronger S2 signal at the same layers as
Llama-3.1-8B-Instruct, (b) a qualitatively different layer-wise profile, or
(c) no meaningful difference?  Relate to the CKA analysis and the
\citet{codaforno2025dual} finding of overlapping subspaces.]

\paragraph{Comparison with CogBias.}
[Discuss: how do our probe results compare with \citet{huang2026cogbias}?
Our probes measure S1/S2 processing mode; theirs target bias susceptibility for
mitigation.  Do the two approaches identify the same layers and directions?
Complementary contributions: they demonstrate practical bias reduction via
steering; we provide a richer mechanistic characterization of the underlying
computation.]

\paragraph{Implications for AI safety.}
[Discuss: understanding when models engage in heuristic versus deliberative
processing has direct implications for trust calibration.  If we can detect
from internal representations whether a model is ``actually reasoning'' on a
given input (vs.\ pattern-matching), this enables more targeted oversight.
However, note the limitation that chain-of-thought faithfulness is low
\citep{jiang2025aha}---the activation-level signal may be more trustworthy than
the token-level signal, but we cannot yet guarantee that high deliberation
intensity entails faithful reasoning.]

\paragraph{Limitations.}
[Discuss: (i)~Model scale: our largest model is 9B parameters.  The
deliberation gradient may manifest differently in models with 70B+ parameters.
(ii)~Benchmark scope: 7 task categories and 284 items cover a limited range of
cognitive biases.  (iii)~SAE coverage: pre-trained SAEs may not capture all
relevant features, especially for the distilled models.  (iv)~Causal evidence:
activation steering demonstrates causal relevance of directions, not circuits.
(v)~Anthropomorphism risk: our ``deliberation intensity'' language is
operational, not a claim about phenomenology.]

\paragraph{Future work.}
[Discuss: (i)~Scaling: apply the framework to 70B+ models.
(ii)~Broader tasks: extend beyond cognitive biases to other domains where
heuristic/deliberative processing differs (e.g., moral reasoning, scientific
reasoning).  (iii)~Circuit-level analysis: move from feature-level to
circuit-level understanding of the deliberation gradient.  (iv)~Temporal
dynamics: finer-grained analysis of how deliberation intensity evolves across
the thinking trace.  (v)~Cross-species comparison: compare the deliberation
gradient in LLMs to neural signatures of dual-process cognition in humans
\citep{deneys2023dual}.]


% ============================================================
% Acknowledgments
% ============================================================

\section*{Acknowledgments}

We thank the Harvard Undergraduate AI Safety (HUSAI) group for supporting
this project.  [Additional acknowledgments: compute sponsors, advisors,
feedback from \ldots.]  Compute was provided by [FASRC / RunPod / \ldots].

% ============================================================
% References
% ============================================================

\bibliographystyle{plainnat}
\bibliography{references}

% ============================================================
% Appendix
% ============================================================

\newpage
% ==============================================================================
% Appendix
% ==============================================================================

\appendix

\section{Benchmark Details}
\label{app:benchmark}

\subsection{Category Descriptions and Examples}

For each of the 7 task categories, we provide the full specification: the
cognitive bias being tested, the conflict/no-conflict manipulation, the
expected S1-lure answer, and a worked example.

\paragraph{CRT variants.}
\emph{Conflict}: ``A notebook and a pen cost \$2.20 together.  The notebook
costs \$2.00 more than the pen.  How much does the pen cost?''
\textbf{S1-lure}: \$0.20.  \textbf{S2-correct}: \$0.10.
\emph{No-conflict control}: ``A notebook and a pen cost \$2.10 together.
The notebook costs \$2.00.  How much does the pen cost?''
\textbf{Answer}: \$0.10.

\paragraph{Base rate neglect.}
\emph{Conflict}: ``In a study of 1000 people, 5 are engineers and 995 are
lawyers.  Tom enjoys math puzzles and building circuits.  What is the
probability that Tom is an engineer?''
\textbf{S1-lure}: high ($>$50\%).  \textbf{S2-correct}: $\approx$0.5\%
(base rate dominates).
\emph{No-conflict control}: ``In a study of 1000 people, 500 are engineers
and 500 are lawyers.  Tom enjoys math puzzles and building circuits.  What
is the probability that Tom is an engineer?''
\textbf{Answer}: depends on diagnosticity, but base rate and description
are congruent.

\paragraph{Belief-bias syllogisms.}
\emph{Conflict (valid--unbelievable)}: ``All flowers need water.  Roses need
water.  Therefore, roses are flowers.''  The conclusion is believable but
the syllogism is \emph{invalid} (affirming the consequent).
\textbf{S1-lure}: valid.  \textbf{S2-correct}: invalid.
\emph{No-conflict (valid--believable)}: ``All mammals are warm-blooded.  All
dogs are mammals.  Therefore, all dogs are warm-blooded.''
\textbf{Answer}: valid.

Table~\ref{tab:category_detail} summarizes the full category distribution
with difficulty breakdowns.

\begin{table}[h]
\centering
\caption{Benchmark category distribution.  Difficulty is rated on a 1--5
  scale; the table shows the number of items at each difficulty level within
  each category.}
\label{tab:category_detail}
\vspace{4pt}
\small
\begin{tabular}{lcccccccc}
\toprule
\textbf{Category} & \textbf{Conflict} & \textbf{Control} & \textbf{Total}
  & \textbf{D1} & \textbf{D2} & \textbf{D3} & \textbf{D4} & \textbf{D5} \\
\midrule
CRT variants       & 30 & 30 & 60  &  8 & 12 & 16 & 14 & 10 \\
Arithmetic          & 25 & 25 & 50  & 10 & 10 & 10 & 10 & 10 \\
Syllogisms          & 25 & 25 & 50  & 10 & 10 & 15 & 10 &  5 \\
Base rate neglect   & 20 & 20 & 40  &  6 &  8 & 12 &  8 &  6 \\
Anchoring           & 15 & 15 & 30  &  6 &  6 &  8 &  6 &  4 \\
Framing effects     & 15 & 15 & 30  &  6 &  6 &  8 &  6 &  4 \\
Conjunction fallacy & 12 & 12 & 24  &  4 &  6 &  6 &  4 &  4 \\
\midrule
\textbf{Total}      & \textbf{142} & \textbf{142} & \textbf{284}
  & 50 & 58 & 75 & 58 & 43 \\
\bottomrule
\end{tabular}
\end{table}

\subsection{Benchmark Construction Details}

Items were generated via a controlled template perturbation process:
(i)~identify the abstract mathematical or logical structure of each
classic problem (e.g., ``$X + Y = T$; $X - Y = D$'');
(ii)~instantiate with novel surface features (objects, names, domains)
drawn from a curated list of 200 nouns and 50 proper names to avoid
thematic overlap;
(iii)~for base rate neglect items, verify all posterior probabilities
analytically using Bayes' theorem with the specified priors and
likelihoods;
(iv)~generate the matched no-conflict control by removing or neutralizing
the heuristic lure (e.g., aligning the base rate with the description, or
presenting a straightforward arithmetic version).

\paragraph{Scoring rubric.}
We use a three-tier scoring system:
\begin{itemize}[leftmargin=*,itemsep=1pt]
  \item \textbf{Binary}: correct (1) or incorrect (0), assessed via
  regex matching against the \texttt{correct\_answer} field.
  \item \textbf{Categorical}: correct / S1-lure / other-wrong / refusal.
  The S1-lure label is assigned when the model's response matches the
  pre-specified \texttt{lure\_answer} regex.  Refusals are detected via
  a keyword list (``I cannot,'' ``I'm not sure,'' etc.).
  \item \textbf{Continuous}: $|r - c| / |l - c|$ where $r$ is the model's
  numerical response, $c$ is the correct answer, and $l$ is the lure
  answer.  Values $> 1$ indicate the response is farther from correct than
  the lure; values near $0$ indicate correctness.  Applicable to numerical
  categories only (CRT, arithmetic, base rate).
\end{itemize}

\subsection{Paraphrase Generation}

Each item has 2--5 paraphrases (mean 3.9 per item, 558 total prompts).
Paraphrases were generated by: (i)~synonym substitution for non-critical
nouns and verbs; (ii)~reordering of premise clauses where logically
equivalent; (iii)~varying the question format (``How much\ldots'' vs.\
``What is the cost of\ldots'').  Critically, paraphrases preserve the
\emph{conflict structure}---the same heuristic lure is present in all
paraphrases of a conflict item.

We verify prompt robustness by computing the intra-class correlation (ICC)
of model responses across paraphrases of the same item.  An ICC $\geq 0.7$
indicates that inter-item variance dominates over inter-paraphrase variance,
confirming that our results are not driven by surface phrasing.

% -------------------------------------------------------------------
\section{Probe Architecture Details}
\label{app:probes}

\subsection{Logistic Regression Probe}

We use \texttt{sklearn.linear\_model.LogisticRegressionCV} with the
following configuration:
\begin{itemize}[leftmargin=*,itemsep=1pt]
  \item \textbf{Penalty}: $\ell_2$ regularization.
  \item \textbf{C range}: 10 values log-spaced from $10^{-4}$ to $10^{4}$,
  selected via nested 5-fold CV (inner loop).
  \item \textbf{Solver}: \texttt{lbfgs} (supports $\ell_2$, efficient for
  $d = 50$ after PCA or $d = 4096$ raw).
  \item \textbf{Max iterations}: 5{,}000.
  \item \textbf{Class weights}: \texttt{balanced} (inversely proportional
  to class frequency; our design is balanced, so this is a no-op for the
  primary task-type target but important for correctness and
  bias-susceptibility targets where classes may be imbalanced).
  \item \textbf{Stratification}: 5-fold, stratified by the Cantor pairing
  of (target label, task category) to ensure each fold has proportional
  representation of all category $\times$ label combinations.
\end{itemize}

\subsection{Mass-Mean Baseline}

The mass-mean probe computes class centroids:
\begin{equation}
  \bar{\mathbf{x}}_c = \frac{1}{|S_c|} \sum_{i \in S_c} \mathbf{x}_i,
  \quad c \in \{\text{conflict}, \text{no-conflict}\},
\end{equation}
and assigns each test activation to the class of the nearest centroid in
cosine distance.  This is a non-parametric baseline that requires no
training and has no hyperparameters.

\subsection{MLP Probe}

\begin{itemize}[leftmargin=*,itemsep=1pt]
  \item \textbf{Architecture}: Input $\to$ Linear(256) $\to$ ReLU $\to$
  Dropout(0.1) $\to$ Linear(1) $\to$ Sigmoid.
  \item \textbf{Optimizer}: Adam, learning rate $10^{-3}$, weight decay
  $10^{-4}$.
  \item \textbf{Training}: 200 epochs maximum, batch size 64, early
  stopping with patience 15 on validation loss (using the same CV folds
  as the logistic probe).
  \item \textbf{Loss}: binary cross-entropy with class-balanced weights.
\end{itemize}

\subsection{Contrast-Consistent Search (CCS)}

CCS \citep{burns2022discovering} finds a linear direction in activation
space that maximizes contrast consistency between paired items.  For each
matched pair $(i, i')$ with activations $(\mathbf{x}_i, \mathbf{x}_{i'})$,
CCS projects onto a direction $\mathbf{w}$ and applies a sigmoid
$p_i = \sigma(\mathbf{w}^\top \mathbf{x}_i + b)$.  The CCS loss is:
\begin{equation}
  \mathcal{L}_{\text{CCS}} = \mathbb{E}\!\left[\min(p_i, 1 - p_i)^2\right]
  + \mathbb{E}\!\left[(p_i + p_{i'} - 1)^2\right].
\end{equation}
The first term encourages confident predictions; the second enforces that
paired items receive complementary probabilities.  We optimize with 10
random restarts (different $\mathbf{w}$ initializations) and select the
direction achieving the lowest loss.

\subsection{Per-Layer Probe Accuracy}

Full tables of ROC-AUC at every layer $\times$ position $\times$ model
$\times$ target combination, including selectivity values and
leave-one-category-out results, are provided in the supplementary data
release.

% -------------------------------------------------------------------
\section{SAE Feature Details}
\label{app:sae}

\subsection{Pre-Trained SAE Sources}

\begin{table}[h]
\centering
\caption{Pre-trained sparse autoencoders used in this study.}
\label{tab:sae_sources}
\vspace{4pt}
\small
\begin{tabular}{llll}
\toprule
\textbf{SAE} & \textbf{HuggingFace ID} & \textbf{Target model} & \textbf{Width} \\
\midrule
Llama Scope & \texttt{fnlp/Llama-3\_1-8B-Base-LXR-32x} & Llama-3.1-8B-Base & 32$\times$ ($\sim$131K) \\
Gemma Scope & \texttt{google/gemma-scope-9b-it-res} & Gemma-2-9B-IT & 16K--1M \\
\bottomrule
\end{tabular}
\vspace{2pt}
\raggedright
\footnotesize{Llama Scope was trained on the base model, not Instruct.
We apply it to both Llama-3.1-8B-Instruct and R1-Distill-Llama-8B after
verifying reconstruction fidelity (Section~\ref{sec:sae}).  For
R1-Distill-Qwen-7B, no suitable pre-trained SAE exists; this model is
analyzed via probes and attention only.}
\end{table}

\subsection{Reconstruction Fidelity Check}

For each SAE--model pair, we verify reconstruction quality on a random
sample of 500 activations per layer.  Two criteria must be met:
\begin{itemize}[leftmargin=*,itemsep=1pt]
  \item \textbf{Explained variance}: $R^2 = 1 - \|\mathbf{x} -
  \hat{\mathbf{x}}\|^2 / \|\mathbf{x} - \bar{\mathbf{x}}\|^2 \geq 0.50$.
  \item \textbf{Normalized MSE}: $\text{NMSE} = \|\mathbf{x} -
  \hat{\mathbf{x}}\|^2 / \|\mathbf{x}\|^2 \leq 2\times$ the NMSE
  reported in the original SAE publication on in-distribution data.
\end{itemize}
Layers failing either criterion are excluded from SAE-based claims and
flagged in the results tables.

\subsection{Falsification Protocol Details}

Following \citet{ma2026falsification}, we apply a token-injection
falsification test to every SAE feature that survives differential
activation testing (BH-FDR $q < 0.05$, $|r_{\text{rb}}| \geq 0.3$):
\begin{enumerate}[leftmargin=*,itemsep=1pt]
  \item Identify the top-3 activating tokens for the candidate feature
  across all benchmark items.
  \item Construct 100 random non-cognitive-bias texts (drawn from
  OpenWebText) of length 128--512 tokens.
  \item Inject each of the top-3 tokens at 5 random positions within each
  text (total: $3 \times 100 \times 5 = 1{,}500$ injections per feature).
  \item Measure the feature's activation at each injection site.
  \item If the feature activates at $\geq 50$\% of its mean activation on
  benchmark items in $\geq 50$\% of injections, classify it as
  \emph{falsified} (token-level artifact).
\end{enumerate}

\noindent
The feature survival pipeline is:
\begin{center}
\small
All features $\xrightarrow{\text{Mann--Whitney } U}$ significant (BH-FDR)
$\xrightarrow{|r_{\text{rb}}| \geq 0.3}$ large effect
$\xrightarrow{\text{falsification}}$ genuine processing-mode features.
\end{center}

\subsection{Top-$K$ Feature Lists}

For each model and layer: the top 20 differentially activated features,
their activation statistics, top-activating tokens, and falsification test
results are provided in the supplementary data release.

\subsection{Cross-Model Feature Matching}

For features that survive falsification in both a Llama Scope and Gemma
Scope model, we compute (i)~the cosine similarity of decoder vectors
(after projecting to a shared space via CCA) and (ii)~the Pearson
correlation of per-item activation patterns on the shared 284-item
benchmark.  Feature pairs with cosine similarity $> 0.5$ \emph{and}
activation correlation $> 0.5$ are classified as matched cross-model
features.

% -------------------------------------------------------------------
\section{Attention Head Classification}
\label{app:attention}

For each model, we provide full tables of per-head Mann--Whitney $U$
statistics, BH-FDR-corrected $p$-values, rank-biserial correlations across
all 5 metrics, and the resulting classification (S2-specialized /
S1-specialized / unspecialized).  Results are reported at both query-head
and KV-group granularity.

For GQA models (Llama-3.1 and R1-Distill-Llama, with 32 query heads and 8
KV groups), KV-group-level statistics are computed by averaging attention
weights across the 4 query heads within each group before computing entropy
and Gini metrics.  This is the conservative granularity used for
confirmatory analysis (H5).

Full per-head classification tables are provided in the supplementary data
release, organized as: model $\times$ layer $\times$ head $\times$
\{metric, $U$-statistic, $p_{\text{raw}}$, $p_{\text{BH}}$,
$r_{\text{rb}}$, classification\}.

% -------------------------------------------------------------------
\section{Representational Geometry Supplement}
\label{app:geometry}

\subsection{PCA and UMAP Visualizations}

We provide 2D PCA and UMAP projections of residual stream activations at
every 4th layer (layers 0, 4, 8, \ldots) for all four models.  Items are
colored by condition (conflict vs.\ no-conflict).  Each PCA/UMAP panel is
accompanied by a random projection baseline: the same data projected to 2D
via a random Gaussian matrix, repeated 100 times.  We display the median
random projection alongside the PCA/UMAP to calibrate visual cluster
expectations.

UMAP hyperparameters: $n_{\text{neighbors}} = 15$, $\text{min\_dist} =
0.1$, metric = cosine.

\subsection{CKA Matrices}

We compute linear CKA \citep{kornblith2019similarity} between
Llama-3.1-8B-Instruct and R1-Distill-Llama-8B on a per-layer basis.
Because both models have 32 layers, this yields a $32 \times 32$ CKA
matrix.  We compute separate matrices for conflict items ($N = 142$) and
no-conflict items ($N = 142$).  The prediction (Section~\ref{sec:geometry})
is that CKA values are lower (more divergent representations) for conflict
items in mid-to-late layers, reflecting reasoning-distillation-specific
changes.

\subsection{Intrinsic Dimensionality}

We estimate the intrinsic dimensionality (ID) of conflict and no-conflict
activation point clouds at each layer using the Two-NN estimator
\citep{facco2017estimating}.  Two-NN estimates ID from the ratio of
distances to the first and second nearest neighbors, requiring no
hyperparameters.  We report point estimates with bootstrap 95\% CIs (1{,}000
resamples).  Lower ID for one condition would indicate more stereotyped
(lower-dimensional) processing; higher ID would suggest diverse processing
strategies.

% -------------------------------------------------------------------
\section{Causal Intervention Details}
\label{app:causal}

\subsection{Dose-Response Curves}

For each non-falsified S2-associated feature, we provide the full
dose-response curve: $P(\text{correct})$ on the 142 conflict items at
each steering coefficient $\alpha \in \{0, 0.5, 1.0, 2.0, 3.0, 5.0\}$.
The random-direction control (mean over 10 random unit vectors) is plotted
alongside.  Error bars denote $\pm 1$ standard deviation across 3 seeds.

\subsection{Side-Effect Analysis}

To check whether S2-feature steering degrades general capabilities, we
evaluate steered models on two held-out benchmarks:
\begin{itemize}[leftmargin=*,itemsep=1pt]
  \item \textbf{MMLU} (5-shot): accuracy on a 500-item subsample spanning
  all 57 categories.
  \item \textbf{HellaSwag} (10-shot): accuracy on a 200-item subsample.
\end{itemize}
A degradation of $>$5 percentage points on either benchmark at the best
steering coefficient is flagged as a meaningful side effect.

\subsection{Per-Category Steering Breakdown}

We report $\Delta P(\text{correct})$ under the best steering coefficient
broken down by all 7 benchmark categories.  This reveals whether causal
effects are category-general (supporting a domain-general S2 mechanism) or
concentrated in specific categories (suggesting category-specific shortcuts).

% -------------------------------------------------------------------
\section{Reproducibility}
\label{app:reproducibility}

\subsection{Compute Budget}

\begin{table}[h]
\centering
\caption{Compute budget for the full experimental pipeline (4 models,
  284 items, 3 seeds where applicable).}
\label{tab:compute}
\small
\begin{tabular}{llcr}
\toprule
\textbf{Stage} & \textbf{Hardware} & \textbf{Time (h)} & \textbf{Cost (\$)} \\
\midrule
Activation caching (4 models) & 1$\times$ A100-80GB & 15--30 & 25--50 \\
Attention extraction (4 models) & 1$\times$ H100-80GB & $\sim$4 & $\sim$10 \\
Probe training (4 models $\times$ 4 targets $\times$ 3 seeds) & CPU (64-core) & $\sim$2 & 0 \\
SAE feature analysis (2 models with SAEs) & 1$\times$ A100-80GB & $\sim$2 & $\sim$3 \\
SAE falsification (top-K features) & 1$\times$ A100-80GB & $\sim$3 & $\sim$5 \\
Causal interventions (5 coefficients $\times$ 3 seeds) & 1$\times$ A100-80GB & 10--20 & 15--30 \\
Geometry analysis (PCA, CKA, Two-NN) & CPU (64-core) & $\sim$4 & 0 \\
\midrule
\textbf{Total} & & $\sim$\textbf{60} & \textbf{80--120} \\
\bottomrule
\end{tabular}
\vspace{2pt}
\raggedright
\footnotesize{GPU costs assume RunPod on-demand pricing (\$1.64/h for
A100-80GB, \$2.49/h for H100-80GB).  CPU stages run on a local workstation
at no incremental cost.}
\end{table}

\paragraph{Storage estimates.}
Residual stream activations for all 4 models at all layers and token
positions occupy approximately 4\,GB in BFloat16 HDF5 format.  Per-head
attention entropy metrics add $\sim$200\,MB.  SAE feature activations
(sparse) add $\sim$500\,MB.  Total storage for the complete experimental
pipeline is $\sim$5\,GB.

\subsection{Hydra Configurations}

All experiments are configured via Hydra YAML files.  The key configuration
groups are:
\begin{itemize}[leftmargin=*,itemsep=1pt]
  \item \texttt{config/model/}: per-model configs specifying HuggingFace
  ID, number of layers, head counts, hidden dimension, and SAE source.
  \item \texttt{config/extract/}: activation extraction settings (token
  positions, BFloat16 precision, HDF5 schema version).
  \item \texttt{config/probe/}: probe hyperparameters (C range, MLP
  architecture, CCS restarts, CV folds, seeds).
  \item \texttt{config/sae/}: SAE analysis settings (reconstruction
  fidelity thresholds, falsification parameters, FDR level).
  \item \texttt{config/attention/}: entropy metrics, GQA handling,
  sliding-window layer separation.
  \item \texttt{config/geometry/}: PCA components, UMAP parameters,
  CKA settings, permutation count.
  \item \texttt{config/causal/}: steering coefficients, random direction
  count, side-effect benchmarks.
\end{itemize}
Full configs are serialized alongside every result file and are included
in the supplementary code release.

\subsection{Random Seeds}

All stochastic components use 3 seeds: \{42, 137, 2024\}.  Seeds are set
globally via \texttt{s1s2.utils.seed.set\_global\_seed()} which seeds
\texttt{torch}, \texttt{numpy}, and \texttt{random}.

\subsection{Software Versions}

Key dependencies (pinned in \texttt{environment.yml}):
\begin{itemize}[leftmargin=*,itemsep=1pt]
  \item Python 3.11
  \item PyTorch $\geq$ 2.3.0 (CUDA 12.1)
  \item Transformers $\geq$ 4.44.0
  \item SAELens $\geq$ 0.4.0
  \item scikit-learn $\geq$ 1.5.0
  \item NumPy $\geq$ 1.26.0
  \item h5py $\geq$ 3.11.0
  \item Weights \& Biases $\geq$ 0.17.0
  \item Hydra-core $\geq$ 1.3.2
\end{itemize}
The exact \texttt{environment.yml} is included in the code release and can
reproduce the full conda environment.


\end{document}
